\input{JMP_slides/preamble.tex}
\title{Fertility and Housing Tenure Choice}
\author{Tommaso De Santo}
\date{Draft, October 2026}
\begin{document}
{
\setbeamertemplate{footline}{}
\maketitle
}

\begin{frame}{Housing and Fertility}
\begin{center}
\includegraphics[width=0.45\textwidth]{../../Latex/newsweek.png}\hspace{-1em}%
\includegraphics[width=0.45\textwidth]{../../Latex/housingwire.png}\\[-0.8em]
\includegraphics[width=0.45\textwidth]{../../Latex/vox.png}\hspace{-1em}%
\includegraphics[width=0.45\textwidth]{../../Latex/cityjournal.png}
\includegraphics[width=0.8\textwidth]{../../Latex/times.png}
\end{center}
\end{frame}

\begin{frame}[label=intro-channels]{Housing, Fertility, and Tenure Choice}
\textbf{Housing costs affect fertility}, but tenure and wealth matter:
\begin{itemize}
  \item Earlier access to homeownership increases family size \cit{(Fazio et al.\ 2025)}
  \item Credit access raises buying, births, and home size \cit{(Hacamo 2021, Couillard 2025)}
\end{itemize}
\medskip
Children increase how big of a house you need. Key channels:
\begin{itemize}
  \item Owning: tenure segmentation (big houses are often only for sale) \hyperlink{big-units-owned}{\beamergotobutton{Data}}
  \item Wealth: borrowing constraints affect purchase ability
  \item Intergenerational distribution: old people own most of the large dwellings \hyperlink{large-homes-data}{\beamergotobutton{Data}}
\end{itemize}
\vfill
How much can physical space affect fertility?\\
What kind of housing policy is good for natality?
\end{frame}

\begin{frame}{This Paper}
\small
\alert{Question:} how much does access to space shape fertility, and can housing policy change it?

\medskip
\alert{Model:} a lifecycle economy with one housing market, in which children need space
\begin{itemize}\setlength\itemsep{2pt}
  \item Households choose when to try for children; children raise the space they need.
  \item Renting is capped in size: large homes must be bought (tenure segmentation).
  \item Buying requires a down payment: wealth and its timing matter (borrowing constraints).
  \item Households age, save, and bequeath: large homes end up with the old (intergenerational allocation).
\end{itemize}

\medskip
\alert{Quantification:} the U.S. as a transition to a smaller population
\begin{itemize}\setlength\itemsep{2pt}
  \item Calibrate the 2007 economy; fit the 2007--2023 fertility decline with preference shocks.
  \item Treat 2023 as a point on the transition and project forward.
  \item Policy: can a property tax that shifts space to young families raise births?
\end{itemize}
\end{frame}

% ======== v8: placeholder for the new framing (open) ========
\begin{frame}{The Argument}
\hole{Framing being worked out (family-sized rentals separate having children from holding housing equity). This frame and the
``This Paper'' frame will be rewritten once the argument is settled; nothing here is final.}
\vfill
\begin{center}\placeholder{[Placeholder for the argument]}\end{center}
\vfill
\end{frame}

\section{Model}

\begin{frame}[plain,c,noframenumbering]{}
\centering
{\Large\textbf{Model}}
\end{frame}

\begin{frame}{Environment}
\begin{itemize}
  \item Discrete time, overlapping generations of households
  \item Adult age $a\in[0,J]$ with survival probability $s_a$; working age $a\in[0,a_R]$; fertile window $a\in[a_F^0,a_F^1]$
  \item Children:
  \begin{itemize}
    \item mature stochastically; housing need depends on children at home
    \item are dependents and do not work
  \end{itemize}
  \item Households differ in financial wealth, inherited housing product, family composition, and persistent idiosyncratic labor efficiency $z_t$
  \item Households choose:
  \begin{itemize}
    \item \textbf{Fertility}: whether to attempt one additional birth, subject to EV shocks
    \item \textbf{Tenure}: rent or choose a product from a discrete owner menu, subject to EV shocks
    \item \textbf{Savings}: next-period liquid wealth $b_{t+1}$; mortgages secured by the house
  \end{itemize}
\end{itemize}
\end{frame}

% Adopted stationary anchor: production/reference_inputs/bundle.json (sequential_births, n_parity, independent child states).
% Bequest convention: engine/household.py::bequest_utility_vec; v8 (Oct 7): Estate A (net of selling cost) is the current production contract.
\begin{frame}{Preferences}
\footnotesize
\setlength{\abovedisplayskip}{2pt}\setlength{\belowdisplayskip}{2pt}
$n_t$: children ever born; $m_t$: children currently at home.

\smallskip
\textbf{Flow utility.} The consumption--space composite is
\[
\mathcal C_t=\frac{c_t^{\alpha_0}(s_t^{\mathrm{net}})^{1-\alpha_0}}{e(m_t)},\qquad
u_t=\frac{\mathcal C_t^{1-\sigma}}{1-\sigma}+\psi_t m_t^{1-\gamma},\quad \gamma\in[0,1).
\]
\begin{itemize}\setlength\itemsep{1pt}
  \item $\alpha_0$: consumption share; $\sigma$: utility curvature; $\gamma$: curvature of the child benefit.
  \item $e(m_t)=((2+\omega m_t)/2)^\nu$; $\omega$ weights children and $\nu$ governs household economies of scale.
  \item $\psi_t m_t^{1-\gamma}$ is the direct benefit; a successful first birth carries the one-time utility cost $\xi$.
\end{itemize}

\smallskip
\textbf{Housing services.} With children at home, the floor is $h_P$; $\chi$ scales owners' residual space:
\[
s_t^{\mathrm{net}}=\mathbf{1}\{\mathrm{own}\}\chi\bigl[h_t-h_P\mathbf{1}\{m_t>0\}\bigr]
+\mathbf{1}\{\mathrm{rent}\}\bigl[h_t^R-h_P\mathbf{1}\{m_t>0\}\bigr]>0,\qquad \chi>0.
\]

\smallskip
\textbf{Bequests.} Over the estate net of selling costs, $w^e_{t+1}=b_{t+1}+(1-\psi^s)P_t h_t$, normalized to zero at a zero estate:
\[
B(w^e_{t+1})=\theta_0\frac{(\theta_1+\max\{w^e_{t+1},0\})^{1-\sigma}-\theta_1^{1-\sigma}}{1-\sigma}.
\]
\end{frame}

% Adopted stationary anchor: engine/household.py (_build_housing_stage_ctx, native purchase-income map, native_due_owner_floor); bundle.json (R_gross, phi, psi, unsecured_credit_limit).
\begin{frame}[label=housing-constraints]{Housing, Tenure, and Budget Constraints}
\titlelink{big-units-owned}{Data}
\small
\setlength{\abovedisplayskip}{2pt}\setlength{\belowdisplayskip}{2pt}
\textbf{Income and saving.}
\begin{itemize}\setlength\itemsep{1pt}
  \item After-tax earnings $y_t^d(a,z_t)$ and a lump-sum transfer $T_t$.
  \item One-period bond $b_{t+1}$ at gross rate $R_b$; owners borrow only against the house.
\end{itemize}

\smallskip
\textbf{Owning.} Choose $h_t$ from a discrete owner-product menu, with largest size $\bar h$, arriving with inherited $h_{t-1}$.
\begin{itemize}\setlength\itemsep{1pt}
  \item The house depreciates at rate $\delta_H$ and is taxed at rate $\tau_t^p$.
  \item Selling costs a fraction $\psi^s$ of the sale value; $\phi$ is the financed share.
  \item Let $S_t=(1-\psi^s)P_t h_{t-1}$ in a transaction and $S_t=0$ otherwise. A change to an owner product requires
  \[
  R_b b_t+S_t+y_t^d+T_t\geq(1-\phi)P_t h_t,\qquad
  b_{t+1}\geq-\phi P_t h_t.
  \]
\end{itemize}

\smallskip
\textbf{Renting.} Rent $h_t^R\in[0,\bar h^R]$ at rent $r_t$ per unit, with $\bar h^R<\bar h$.
\begin{itemize}\setlength\itemsep{1pt}
  \item The largest homes are reachable only through ownership.
\end{itemize}

\end{frame}

\begin{frame}{Saving and Sale Solvency}
\small
Renters have no unsecured credit, and selling to rent must settle inherited debt:
\[
b_{t+1}\geq0,\qquad R_b b_t+S_t\geq0\quad\text{when an owner sells to rent}.
\]
Current income can finance a purchase; it cannot cover a shortfall when selling to rent.

\medskip
An owner retaining the same home can roll over debt beyond the collateral limit,
but cannot increase that excess debt:
\[
b_{t+1}\geq\min\{b_t,-\phi P_t h_t\}.
\]
When death is possible, the estate after selling costs must be nonnegative:
\[
b_{t+1}\geq\max\{\min(b_t,-\phi P_t h_t),-(1-\psi^s)P_t h_t\}.
\]
\end{frame}

\begin{frame}{Fertility and Child Aging}
\small
\begin{itemize}\setlength\itemsep{0.6em}
  \item During fertile ages, households choose each period whether to attempt at most one additional birth; children ever born are top-coded at three.
  \item The choice is discrete, with an extreme-value taste shock $\varepsilon_t^F$ of scale $\kappa_1$ for the first birth and $\kappa_C$ for later births.
  \item An attempt succeeds with age-specific probability $\pi_a$ and adds one child.
  \item Children mature stochastically: each dependent child leaves the household with a fixed probability each period, independently of the others.
\end{itemize}
\end{frame}

\begin{frame}{Firms and Earnings}
\small
\begin{itemize}
  \item Competitive firms produce with effective labor:
  \[
  Y_t=A L_t,\qquad w=A,\qquad L_t=\int_{a<a_R} e_a z_t\,dG_t.
  \]
  \item Age- and productivity-dependent earnings:
  \[
  y_t^d(a,z_t)=\begin{cases}
  (1-\tau^{w})\,w e_a z_t,&a<a_R,\\
  \varpi_t,&a\geq a_R.
  \end{cases}
  \]
  \item $e_a$: age--earnings profile; $z_t$: idiosyncratic labor efficiency; $G_t$: distribution of households over states.
  \item Retirees receive a pay-as-you-go pension $\varpi_t$; payroll taxes balance it each period.
\end{itemize}
\end{frame}

\begin{frame}{Housing Supply and Rents}
\small
The stationary supply schedule uses the current user cost:
\[
H^s=H_0\left(\frac{r}{\bar r}\right)^{\eta},\qquad
r=(R_b-1+\delta_H+\tau^p)P.
\]
\begin{itemize}
  \item $H_0$: housing supply at reference rent $\bar r$; $\eta$: supply elasticity.
  \item More generally, a landlord's user-cost identity is
  \[
  r_t=(R_b+\delta_H+\tau_t^p)P_t-P_{t+1}.
  \]
\end{itemize}
\end{frame}

\begin{frame}[t]{%
\only<1>{The Household's Problem: Housing and Savings}%
\only<2>{The Household's Problem: Fertility Choice}%
}
\vspace{0.3em}
\centering
\scalebox{0.85}{%
\begin{tikzpicture}[
  node distance=7mm and 10mm, >=Latex,
  see/.style={
    ellipse, draw, fill=yellow!20,
    minimum height=1.3em, minimum width=1.8cm,
    inner sep=2pt, align=center, font=\normalsize
  },
  choose/.style={
    rectangle, draw, fill=green!15, rounded corners=2pt,
    minimum height=1.5em, minimum width=1cm,
    inner sep=2pt, align=center, font=\normalsize
  },
  line/.style={-Latex, semithick, draw=black!75}
]
\node (fertS) [see] {see $\varepsilon_t^F$};
\node (fertD) [choose, right=of fertS, alt=<2>{draw=red, ultra thick}{}] {choose attempt\\($\pi^F_{t,a}$, $V_t$)};
\node (locS)  [see,    right=of fertD] {birth; see $\varepsilon_t^H$};
\node (locD)  [choose, right=of locS,  alt=<1>{draw=red, ultra thick}{}] {choose $h_t,h_t^R,b_{t+1}$\\($\pi^H_{t,a}$, $V^H_t$)};

\draw[line] (fertS) -- (fertD);
\draw[line] (fertD) -- (locS);
\draw[line] (locS)  -- (locD);
\end{tikzpicture}%
}

\par\vspace{0.35em}
\raggedright
\setlength{\abovedisplayskip}{4pt}\setlength{\belowdisplayskip}{4pt}

\only<1>{%
{\footnotesize
\setlength{\abovedisplayskip}{9pt}\setlength{\belowdisplayskip}{9pt}

\textbf{Conditional on the realized family size, choose housing and savings:}
\[
\begin{aligned}
V^H_t(b_t,h_{t-1},z_t,a,n_t,m_t)
={}&\E_{\varepsilon_t^H}\max_{h_t,h_t^R,b_{t+1}}
\Bigl\{u_t\big(c_t,s_t^{\mathrm{net}};m_t\big)+\varepsilon_t^H(h_t)\\
&+\beta\E_t\!\left[s_a V_{t+1}(b_{t+1},h_t,z_{t+1},a+1,n_{t+1},m_{t+1})
+(1-s_a)B(w^e_{t+1})\right]\Bigr\}.
\end{aligned}
\]

subject to
\[
c_t+b_{t+1}+r_t h_t^R+(\delta_H+\tau_t^p)P_t h_t
=R_b b_t+y_t^d(a,z_t)+T_t+S_t-\mathbf{1}\{h_t\neq h_{t-1}\}P_t h_t,
\]
\medskip
\[
\text{with the financing, solvency, and housing-service constraints above.}
\]
$\varepsilon_t^H$: extreme-value tastes over renting and the owner products, scale $\kappa_T$.
}
}

\only<2>{%
{\footnotesize
\setlength{\abovedisplayskip}{6pt}\setlength{\belowdisplayskip}{6pt}

Waiting and attempting have values
\[
\begin{aligned}
W_t^0&=V^H_t(b_t,h_{t-1},z_t,a,n_t,m_t),\\
W_t^A&=\pi_a\!\left[V^H_t(b_t,h_{t-1},z_t,a,n_t+1,m_t+1)-\xi\1\{n_t=0\}\right]\\
&\quad +(1-\pi_a)V^H_t(b_t,h_{t-1},z_t,a,n_t,m_t).
\end{aligned}
\]
Before fertility tastes, household value is
\[
V_t(b_t,h_{t-1},z_t,a,n_t,m_t)
=\E_{\varepsilon_t^F}\max\{W_t^0+\varepsilon_{t,0}^F,\ W_t^A+\varepsilon_{t,A}^F\}.
\]
Choice probabilities:
\[
\pi^F_{t,a}(b_t,h_{t-1},z_t,n_t,m_t)
=\frac{\exp\{W_t^A/\kappa(n_t)\}}
{\exp\{W_t^A/\kappa(n_t)\}+\exp\{W_t^0/\kappa(n_t)\}},
\]
Here $\kappa(n_t)=\kappa_1$ if $n_t=0$ and $\kappa_C$ otherwise. Expected births equal $\pi_a\pi^F_{t,a}$.
}
}
\end{frame}

\begin{frame}{Equilibrium}
\small
\setlength{\abovedisplayskip}{3pt}\setlength{\belowdisplayskip}{3pt}
Given parameters, an initial population $N_0$, and an initial household distribution $G_0$, an equilibrium is a sequence
\[
\{V_t,\;g_t,\;G_t,\;N_t,\;P_t,\;r_t,\;T_t,\;\varpi_t\}_{t\ge0}
\]
such that, for every $t$:
\begin{enumerate}\setlength\itemsep{3pt}
  \item \textbf{Households.} Given $\{P_s,r_s,T_s,\varpi_s\}_{s\ge t}$, $V_t$ solves the household problem with continuation $V_{t+1}$; the policy functions $g_t$ (fertility, housing, savings) attain it.
  \item \textbf{Markets and government.} Rent is the user cost of housing, the housing market clears, pensions balance payroll taxes, and property-tax revenue is rebated equally.
  \item \textbf{Population.} Births $B_t=\int\pi_a\pi^F_{t,a}\,dG_t$ enter at age zero and persons age with survival $s_a$: $N_{0,t+1}=B_t$, $N_{a+1,t+1}=s_a N_{a,t}$.
  \item \textbf{Consistency.} $G_{t+1}$ follows from $G_t$ under $g_t$ and the idiosyncratic shocks, with entrants making the number of households at each age consistent with $N_{t+1}$.
\end{enumerate}
\hole{Entry: the Oct 3 stationary frame has births replacing 2.1 adult entrants through an entry queue (half at 16, half at 20) and no rebate; the current contract rebates the tax. Reconcile the entry statement with the engine.}
\end{frame}

\section{Quantification}

\begin{frame}[plain,c,noframenumbering]{}
\centering
{\Large\textbf{Quantification}}
\end{frame}


% Initial tables and fit coordinates: verified_final/selected_target_fit.csv and selected_parameters.csv (2026-09-12).
% Fingerprint c0e266d3a0d430343c469d780d1aedb45fa87f8763c9c938889e0c37daa31de2.

\begin{frame}[t]{Fertility in the United States}
\centering
% Display copies from the verified WDI/CPS series, plotted from 1990; the completed-fertility continuation is rescaled to end at the constant period TFR (top-code filled by author).
\alt<2>{\includegraphics[width=\textwidth,height=0.8\textheight,keepaspectratio]{figures/fertility_introduction_projection_1990.pdf}}{\includegraphics[width=\textwidth,height=0.8\textheight,keepaspectratio]{figures/fertility_introduction_history_1990.pdf}}
\end{frame}

\begin{frame}[t]{Initial Steady State}
\begin{minipage}[t][0.06\textheight][t]{\textwidth}
\vspace{0pt}
For intuition, summarize housing costs by $p$ and fertility by $n(p;\psi)$.
\end{minipage}\par
\begin{minipage}[t][0.56\textheight][t]{\textwidth}
\vspace{0pt}\centering
\includegraphics[width=\linewidth,height=0.56\textheight,keepaspectratio]{figures/september_housing_fertility_stage_initial.pdf}
\end{minipage}\par
\begin{minipage}[t][0.16\textheight][t]{\textwidth}
\vspace{0pt}\raggedright\small
At $A$, housing clears at $(S_0,p_0)$ and fertility is at the replacement rate $\bar n$.
\end{minipage}
\end{frame}

\begin{frame}[t,noframenumbering]{Impact}
\begin{minipage}[t][0.06\textheight][t]{\textwidth}
\vspace{0pt}
\mbox{}
\end{minipage}\par
\begin{minipage}[t][0.56\textheight][t]{\textwidth}
\vspace{0pt}\centering
\includegraphics[width=\linewidth,height=0.56\textheight,keepaspectratio]{figures/september_housing_fertility_stage_impact.pdf}
\end{minipage}\par
\begin{minipage}[t][0.16\textheight][t]{\textwidth}
\vspace{0pt}\raggedright\small
A fall from $\psi_0$ to $\psi_1$ lowers fertility at unchanged housing cost: $A\to A'$. Housing costs then adjust at inherited population $S_0$: $A'\to B$.
\end{minipage}
\end{frame}

\begin{frame}[t,noframenumbering]{Demographic Adjustment}
\begin{minipage}[t][0.06\textheight][t]{\textwidth}
\vspace{0pt}
\mbox{}
\end{minipage}\par
\begin{minipage}[t][0.56\textheight][t]{\textwidth}
\vspace{0pt}\centering
\includegraphics[width=\linewidth,height=0.56\textheight,keepaspectratio]{figures/september_housing_fertility_stage_adjustment.pdf}
\end{minipage}\par
\begin{minipage}[t][0.16\textheight][t]{\textwidth}
\vspace{0pt}\raggedright\small
Smaller entering cohorts reduce housing demand. Lower housing costs partly offset the preference decline; $C$ marks the new stationary outcome.
\end{minipage}
\end{frame}

\begin{frame}[label=calibration-overview]{Calibration and Historical Transition}
\titlelink{empirics}{Empirical evidence}
\begin{enumerate}\setlength\itemsep{1.1em}
  \item \textbf{Calibrate the 2007 steady state.}\\
  Match fertility, housing, and wealth moments to determine household parameters.
  \item \textbf{Estimate the sequence of fertility shocks.}\\
  Hold those parameters fixed and match birth rates through 2023 with unanticipated changes in fertility preferences. Carry households and population forward between shocks.
  \item \textbf{Project the economy forward.}\\
  After 2023, hold fertility preferences fixed and solve for prices, housing choices, and population. Compare housing policies from the same 2023 economy.
\end{enumerate}
\medskip
A model period is four calendar years.

\end{frame}

\begin{frame}{Utility Specification}
\small
\[
u_t(c_t,s_t;m_t)=
\frac{\left[\dfrac{c_t^{\alpha_0}\bigl(\chi_t[h_t-h_P\mathbf{1}\{m_t>0\}]\bigr)^{1-\alpha_0}}{e(m_t)}\right]^{1-\sigma}}{1-\sigma}
+\psi_tm_t^{1-\gamma}.
\]
\[
e(m_t)=\left(\frac{2+\omega\, m_t}{2}\right)^{\nu}.
\]
\smallskip
$h_t$: rooms occupied; $\chi_t=\chi$ for owners and $1$ for renters.
Equivalence scale $e(m_t)$ as in Scholz, Seshadri \& Khitatrakun (2006).
\end{frame}


\begin{frame}{Identification}
\footnotesize
Ten parameters and the supply scale $H_0$ are chosen jointly to match ten moments and the completed-fertility normalization; three further moments are validation only.

\smallskip
\textbf{Fertility}
\vspace{-3pt}\begin{itemize}\setlength\itemsep{0pt}
  \item $\psi_0$, preference for children $\to$ completed fertility of 2.1
  \item $\xi$, first-birth cost $\to$ childlessness
  \item $\kappa_1$, first-birth taste dispersion $\to$ mean age at first birth, children ever born by 25
  \item $\kappa_C$, later-birth taste dispersion, and $\gamma$, curvature of the benefit of children $\to$ one-child families
\end{itemize}
\vspace{-2pt}
\textbf{Housing}
\vspace{-3pt}\begin{itemize}\setlength\itemsep{0pt}
  \item $h_P$, space requirement of parents $\to$ rooms added at first birth
  \item $\chi$, owner premium $\to$ ownership rate
  \item $\kappa_T$, tenure taste dispersion $\to$ ownership gap of recent parents
  \item $H_0$, supply scale $\to$ mean rooms
\end{itemize}
\vspace{-2pt}
\textbf{Saving and bequests}
\vspace{-3pt}\begin{itemize}\setlength\itemsep{0pt}
  \item $\beta$, discount factor $\to$ wealth over earnings
  \item $\theta_0$, bequest scale $\to$ bequest flow ($\theta_1$ fixed externally)
\end{itemize}
\hole{Moment-to-parameter pairing is the intended one; informative rank is not yet checked.}
\end{frame}

\begin{frame}{Calibration: Targets and Model}
\footnotesize
\linespread{1}\selectfont
\renewcommand{\arraystretch}{1.12}
\renewcommand{\arraystretch}{0.9}
\input{\tabdirv/fit.tex}\par
\smallskip
\tiny $^{v}$ Validation moment, not targeted. Weighted loss \baseloss; full table with weights in the appendix.\\
Sources: CPS June 2004/2006 (children ever born, women 40--44 and at 25); NCHS natality 2003--2006 (first-birth timing); AHS 2007 (rooms);
ACS 2005/2006 (ownership, recent parents, family rooms); PSID 2005/2007 (wealth over earnings); SCF 2007 (bequests);
PSID event study, years $+3/+4$ versus $-3/-2$ (rooms added at first birth).
\end{frame}
\begin{frame}{Internally Calibrated Parameters}
\footnotesize

\renewcommand{\arraystretch}{1.08}
\input{\tabdirv/params.tex}
\par\smallskip
\scriptsize $^{*}$ Flagged near a search bound in the saved parameter table. Optimization convergence is not certified.
\end{frame}

\begin{frame}{Externally Calibrated Parameters}
\footnotesize
\renewcommand{\arraystretch}{1.0}
\input{\tabdirv/external.tex}
\end{frame}

% ======== October 2026: household choices (the full set of new results is in JMP_slides_new_results.tex) ========
\section{How Households Decide}
\begin{frame}[plain,c,noframenumbering]{}
\centering
{\Large\textbf{How Households Decide}}
\end{frame}

% Source: output/model/jmp_slides_policy_functions_20261007/ (code/model/tools/build_jmp_policy_function_figure.py), saved 14.402 solution, no re-solve.
\begin{frame}[t]{Who Has a First Child: Policy Functions}
\linespread{1.0}\footnotesize\selectfont
Probability that a childless renter has a first birth in the period, as chosen by the household.
\begin{center}
\includegraphics[width=\textwidth,height=0.50\textheight,keepaspectratio]{../output/model/jmp_slides_policy_functions_20261007/first_birth_policy.pdf}
\end{center}
\vspace{-0.6em}
\begin{itemize}\setlength\itemsep{1pt}\setlength\topsep{2pt}
\item Births rise steeply with liquid wealth near the borrowing limit and flatten once a household holds a buffer.
\item Higher earnings shift the curve left: the same child is affordable with less saved. Births peak at 26--33.
\end{itemize}
\runlabel{Fixed price, 2007 steady state. Attempt probability times conception probability; earnings states 0.47, 0.78, 1.29 of mean. Right panel at zero liquid wealth.}
\end{frame}


% Source: output/model/jmp_slides_policy_functions_20261007/child_policy.pdf
%   (code/model/tools/build_jmp_child_policy_figure.py --expected), saved 14.402 solution, no re-solve.
\begin{frame}[t]{What a First Child Changes: Policy Functions}
\linespread{1.0}\footnotesize\selectfont
Childless renter, age 26--29, middle earnings: choices this period with no child and with a first child just born.
\smallskip

\begin{columns}[T,onlytextwidth]
\begin{column}{0.64\textwidth}
\includegraphics[width=\textwidth,height=0.66\textheight,keepaspectratio]{../output/model/jmp_slides_policy_functions_20261007/child_policy.pdf}
\end{column}
\begin{column}{0.33\textwidth}
\begin{itemize}\setlength\itemsep{6pt}
\item With little wealth, a child means renting at the 6-room cap: the household delays buying.
\item With a buffer, a child means buying a bigger home.
\item Trying for a child pays only once the household holds a small buffer; below it, waiting is worth more.
\end{itemize}
\end{column}
\end{columns}
\runlabel{Fixed price, 2007 steady state. Rooms and consumption averaged over the tenure choice; earnings state 0.78 of mean. Value of trying minus waiting:
$\pi_a[V^H(\text{child})-\xi-V^H(\text{no child})]$, utility units; its logit over $\kappa_1$ is the attempt probability on the previous slide.}
\end{frame}

% Sources: output/model/reconciliation_14p40_20261007/precautionary/README.md (headline cells, risk A);
%   output/model/rental_menu_precaution_20261007/tables_price_tests.md (Test 1, risk A column, states 3-8). Base 14.402.
\begin{frame}[t]{Income Risk and the Rent of Family Space}
\linespread{1.0}\footnotesize\selectfont
A first child commits the household to more space, the floor $h_P$, for as long as the child is at home. With risky income and no savings,
that commitment is risky. Experiment: lower the standard deviation of income shocks from 0.48 to 0.35, nothing else.
\smallskip

\begin{columns}[T,onlytextwidth]
\begin{column}{0.48\textwidth}
\scriptsize\setlength{\tabcolsep}{3pt}\renewcommand{\arraystretch}{1.05}
\begin{tabular}{@{}lrr@{}}
\toprule
 & Baseline risk & Lower risk \\
\midrule
First-birth attempts, 22--33 & 0.185 & $+39\%$ \\
\quad renters & 0.119 & $+57\%$ \\
\quad owners & 0.357 & $+23\%$ \\
Childless at 45 (\%) & 19.0 & 7.4 \\
\bottomrule
\end{tabular}
\end{column}
\begin{column}{0.48\textwidth}
\scriptsize\setlength{\tabcolsep}{3pt}\renewcommand{\arraystretch}{1.05}
\begin{tabular}{@{}lrrr@{}}
\toprule
House prices and rents & $-10\%$ & base & $+10\%$ \\
\midrule
Effect of lower risk & $+33\%$ & $+39\%$ & $+44\%$ \\
\bottomrule
\end{tabular}
\end{column}
\end{columns}
\smallskip

\footnotesize
\begin{itemize}\setlength\itemsep{1pt}\setlength\topsep{2pt}
\item Lower risk cuts childlessness at 45 from 19.0\% to 7.4\%; renters with no savings respond most.
\item The more family space costs, the more income risk deters a first birth.
\end{itemize}

\runlabel{Fixed price, rebate on and held at its baseline value. Attempts: first-birth attempt probability of childless households at income states 3--8, same states and weights.
Lower risk: same persistence (0.73), income grid renormalized to mean one. Right: prices and rents moved together.}
\end{frame}

% Source: output/model/rental_menu_precaution_20261007/tables_floor_credit.md ("Credit, transfers and price on one calibration"), base 14.402.
\begin{frame}[t,label=loan-not-insurance]{A Loan Is Not Insurance}
\titlelink{ces-appendix}{CES}
\linespread{1.0}\footnotesize\selectfont
Same households and parameters, house prices fixed. Each row changes one thing.
\smallskip

\scriptsize\setlength{\tabcolsep}{4pt}\renewcommand{\arraystretch}{1.05}
\begin{tabularx}{\textwidth}{@{}Xrrrr@{}}
\toprule
 & Births & Completed & Childless & Own \\
 & on impact & fertility & 46--49 (\%) & 18--29 (\%) \\
\midrule
Benchmark & & & 19.0 & 43.8 \\
\addlinespace[2pt]
House prices and rents $+10\%$ & $-4.4\%$ & $-5.3\%$ & 22.3 & 40.9 \\
Financed share $80\%\to95\%$ & $+0.4\%$ & $-1.4\%$ & 20.2 & 69.0 \\
Renter credit, 1 year of earnings & $+5.1\%$ & $+0.7\%$ & 19.0 & 37.1 \\
Renter credit, 2 years of earnings & $+8.5\%$ & $+0.6\%$ & 19.4 & 34.4 \\
\addlinespace[2pt]
Transfer paid in bad years & $+17.3\%$ & $+14.7\%$ & 9.7 & 40.8 \\
Same cost, paid to everyone & $+7.6\%$ & $+6.6\%$ & 14.8 & 44.0 \\
\bottomrule
\end{tabularx}
\smallskip

\footnotesize
\begin{itemize}\setlength\itemsep{1pt}\setlength\topsep{2pt}
\item Mortgage credit moves ownership, not family size. Renter credit moves births earlier, not more.
\item A transfer of the same cost raises completed fertility twice as much when it pays in bad years.
\end{itemize}

\runlabel{Fixed price, property-tax rebate on and held at its baseline value. Impact: births on the baseline distribution of states. Completed fertility: children ever born at 46--49, stationary cohort.
Renter credit: renters may borrow $D$ = 1 or 2 years of median earnings at 22. Bad-year transfer: 25\% of the gap to median earnings when below it; both transfers cost 0.097 per household per period.}
\end{frame}

% Sources: output/model/ltv_branch_values_20261008/tables_branch_values.md (fp, rebate re-balanced, omega-weighted marginal households);
%   output/model/origination_only_mortgage_20261008/tables.md; price row: project file tables_capitalization.md (ordinary +10%, [0,6]). Base 14.402, fixed price.
\begin{frame}[t]{Why Credit Barely Moves Births}
\linespread{1.0}\footnotesize\selectfont
Rent per room equals the owner's user cost, so the space a child needs costs the same rented or owned, at any loan-to-value.
Credit changes when a household buys, not what the child's space costs.
\smallskip

\scriptsize\setlength{\tabcolsep}{5pt}\renewcommand{\arraystretch}{1.05}
\begin{tabularx}{\textwidth}{@{}Xrrrr@{}}
\toprule
 & Value with & Value & Gain from & Births \\
 & a birth & waiting & a birth & on impact \\
\midrule
Financed share $80\%\to95\%$ & $+0.0081$ & $+0.0068$ & $+0.0014$ & $+0.31\%$ \\
\quad renters & $+0.0096$ & $+0.0083$ & $+0.0013$ & \\
\quad owners, 4+ rooms & $+0.0053$ & $+0.0038$ & $+0.0015$ & \\
Financed share $80\%\to100\%$ & $+0.0180$ & $+0.0123$ & $+0.0057$ & $+1.30\%$ \\
\addlinespace[2pt]
House prices and rents $+10\%$ & & & $-0.111^{a}$ & $-4.41\%$ \\
\bottomrule
\end{tabularx}
\smallskip

\footnotesize
\begin{itemize}\setlength\itemsep{1pt}\setlength\topsep{2pt}
\item Credit is worth about the same to a household whether or not it has the child; the birth-specific part is about a sixth of that, against a taste dispersion $\kappa_1=0.12$.
\item Without cash-out refinancing (credit only at purchase), the same policy gives births $+0.02\%$ on impact and completed fertility $-1.7\%$.
\end{itemize}
\runlabel{Fixed price, rebate re-balanced (the previous slide holds the rebate fixed, hence $+0.4\%$ there for the same policy). Changes in value in utility units, averaged over households at the first- or later-birth margin (weights $\propto$ mass $\times\,\pi_a^2p(1-p)/\kappa$). Gain from a birth: value with a birth minus value of waiting. $^{a}$Ages 22--33 only.}
\end{frame}

% v8: H100 two-shock transition at the 14.402 base (output/model/transition_h100_14402_20261007/, Transition thread, Oct 7).
\begin{frame}[label=transition-short]{Fertility Transition: 2007--2063}
\titlelink{algo-ss}{Algorithm}
\centering
\includegraphics[width=\textwidth,height=0.80\textheight,keepaspectratio]{../output/model/transition_h100_14402_20261007/pair_0.1311_0.110/fertility_fit_2007_2063.pdf}
\runlabel{Current base (loss \baseloss), GE, rebate rebalanced every date, 100-period horizon per stage. Shocks fitted to the 2011--14 and 2019--22 period TFR.}
\end{frame}

% Source: output/model/transition_h100_14402_20261007/path_table_0.1311_0.110.md (H100, 14.402 base).
\begin{frame}[t]{Two Preference Shocks: 2007--2067}
\footnotesize
The preference for children falls twice, unanticipated and permanent each time, with perfect foresight in between:
$\psi$ from 0.156 to 0.131 in 2007 and to 0.110 in 2015. Each shock is set to hit one period fertility target.
\smallskip

\scriptsize\setlength{\tabcolsep}{5pt}\renewcommand{\arraystretch}{1.05}
\begin{tabularx}{\textwidth}{@{}Xrrrrrrr@{}}
\toprule
 & 2007 & 2011 & 2015 & 2019 & 2023 & 2043 & 2067 \\
\midrule
Period TFR, model & 1.863 & 1.862 & 1.658 & 1.650 & 1.646 & 1.695 & 1.813 \\
Period TFR, target & & 1.861 & & 1.646 & & & \\
House price (2007 steady state 0.779) & 0.772 & 0.769 & 0.760 & 0.754 & 0.748 & 0.696 & 0.615 \\
Households (2007 = 1) & 1.000 & 1.000 & 1.000 & 1.000 & 0.997 & 0.939 & 0.832 \\
\bottomrule
\end{tabularx}
\smallskip

\footnotesize Fertility falls on impact; the house price falls slowly, as smaller cohorts reach the housing market. The economy reaches its new steady state
(price 0.485, population 0.59 of 2007) by about 2400.

\runlabel{Current base (loss \baseloss), GE, instantaneous supply, rebate balanced every date, 100-period horizon per stage; stage-2 residuals below $1.1\times10^{-5}$, terminal check passed.}
\end{frame}

\begin{frame}{Full Transition: 2007--2403}
\centering
\includegraphics[width=\textwidth,height=0.80\textheight,keepaspectratio]{../output/model/transition_h100_14402_20261007/pair_0.1311_0.110/full_transition_four_panel.pdf}
\runlabel{Current base (loss \baseloss), GE, rebate rebalanced every date, 100-period horizon per stage.}
\end{frame}

% Model column: output/model/profile_2023_base14402/moments_2023_base14402.csv (H100 baseline replay, 14.402 base). Sept 14 version in git (e49c148c).
\begin{frame}{2023 Data and Model}
\footnotesize
\renewcommand{\arraystretch}{1.05}
\begin{tabularx}{\textwidth}{@{}Xrr@{}}
\toprule
Moment & Data & Model \\
\midrule
Period TFR (fitted) & 1.646 & 1.646 \\
Completed fertility, ages 40--44 (model 42--45) & 1.92 & 1.69 \\
Childless, ages 40--44 (model 42--45) (\%) & 18.8 & 25.3 \\
Mean age at first birth (years) & 28.09 & 26.85 \\
\addlinespace[2pt]
Mean occupied rooms (capped at 9) & 5.58 & 5.54 \\
Ownership, heads 30--55 (model 30--53) (\%) & 58.7 & 69.3 \\
\bottomrule
\end{tabularx}
\runlabel{Current base (loss \baseloss), 2023 on the 100-period two-shock path, GE, rebate balanced every date. Model ages are four-year groups. The one-child share, first births at 30+, room responses and wealth rows of the earlier version are not in this readout.}
\end{frame}

\begin{frame}{Lifecycle Fit in 2023}
\centering
\includegraphics[width=\textwidth,height=0.85\textheight,keepaspectratio,trim=0 25 0 0,clip]{../output/model/profile_2023_base14402/lifecycle_2023.pdf}
\runlabel{Current base, 2023 on the 100-period two-shock path. Children at home: model dependents; ACS resident own children under 18.}
\end{frame}

\begin{frame}{Intergenerational Allocation: Model vs Data}
\centering
% Current-base figure: output/model/profile_2023_base14402/build_figures.py.
\includegraphics[width=0.9\linewidth,height=0.5\textheight,keepaspectratio]{../output/model/profile_2023_base14402/intergenerational_allocation_2023.pdf}
\par\centerline{\scriptsize Share of owner-occupied homes with $6+$ rooms; ACS 2023.}
\vspace{4pt}
\raggedright\footnotesize
\begin{itemize}\setlength\itemsep{1pt}
  \item Large homes sit with older households without children; young families hold few.
  \item Coven et al.\ (2025): a rebated property tax moves space from old owners to young households.
  \item Those young households are the ones deciding on children, and children need space.
\end{itemize}
\runlabel{Model bars: current base, 2023 on the 100-period two-shock path.}
\end{frame}

\begin{frame}{Policy}
\small
\textbf{The experiment of Coven, Golder, Gupta \& Ndiaye (2025)}
\begin{itemize}\setlength\itemsep{2pt}
  \item Tax housing by value; rebate the revenue equally to all households.
  \item Older owners pay more than they get back; young renters get back more than they pay.
  \item House prices fall; constrained young households buy sooner or buy bigger.
\end{itemize}

\medskip
\textbf{This paper}
\begin{itemize}\setlength\itemsep{2pt}
  \item Same experiment: property tax from 1.06\% to 2\%, equal rebate, from the 2023 economy.
  \item Follow prices, rents, who lives in the large homes, and births along the path.
  \item Does the space that moves to young families turn into children?
\end{itemize}
\end{frame}

% Source: output/model/policy_transition_battery_20261008/h100_torch/ (table_h100.md, tax_vs_credit_h100.pdf), base 14.402, from the 100-period 2023 state.
\begin{frame}{Policy Results}
\centering
\includegraphics[width=.95\textwidth,height=.74\textheight,keepaspectratio]{../output/model/policy_transition_battery_20261008/h100_torch/tax_vs_credit_h100.pdf}
\runlabel{Current base (loss \baseloss), GE from the 2023 state, 100 periods, elastic supply, rebate balanced every date. Property tax 1.06\% to 2\%.}
\end{frame}

% Source: output/model/policy_transition_battery_20261008/h16_mac/ (table_h16.md, tax_vs_credit_by_supply_h16.pdf), 16-period runs.
\begin{frame}{Policy Results by Housing Supply \textnormal{\small(16-period runs)}}
\centering
\includegraphics[width=.92\textwidth,height=.74\textheight,keepaspectratio]{../output/model/policy_transition_battery_20261008/h16_mac/tax_vs_credit_by_supply_h16.pdf}
\runlabel{Current base, GE from the 2023 state, 16 periods (dates after about 2047 shaped by the terminal condition); rebate balanced every date. Slow stock: $H_t=(1-d)H_{t-1}+dH_{2023}(q_t/q_{2023})^{0.63}$.
Fixed and slow arms use the inherited-distribution projection; the slow-stock credit run stopped at its 20-map cap (residual $1.5\times10^{-4}$).}
\end{frame}

% Source: \pol macros, generated from output/model/reconciliation_14p40_20261007 (Blocks 2 and 3), base 14.402.
\begin{frame}[t]{Property Tax versus Mortgage Credit in the Long Run}
\linespread{1.0}\footnotesize\selectfont
Stationary equilibrium at the 2023 preference for children, with and without each policy. Births are pinned at replacement, so the response shows in population and prices.
\medskip

\small\setlength{\tabcolsep}{6pt}\renewcommand{\arraystretch}{1.15}
\begin{tabularx}{\textwidth}{@{}Xrrr@{}}
\toprule
 & House price & Population & Own 18--29 \\
\midrule
Property tax $1.06\%\to2\%$, rebated & $\pol{end_ptax2.price_pct.1}\%$ & $\pol{end_ptax2.pop_pct.1}\%$ & $\pol{end_ptax2.own_18_29_pp.1}$ pp \\
Financed share $80\%\to95\%$ & $\pol{end_phi095.price_pct.1}\%$ & $\pol{end_phi095.pop_pct.1}\%$ & $\pol{end_phi095.own_18_29_pp.1}$ pp \\
\bottomrule
\end{tabularx}
\medskip

\footnotesize
\begin{itemize}\setlength\itemsep{2pt}
\item At fixed prices the tax lowers births ($\pol{ptax2_fp.births_pct}\%$, 2007 economy): the whole tax passes into rents.
In equilibrium the house price falls, and the tax supports a larger population.
\item Credit raises young ownership by 26 points and supports a smaller population.
\end{itemize}

\runlabel{General equilibrium: house price, pensions and the rebate solved; housing supply responds to price. 2023 preference for children $\psi=0.110$ from the two-shock fit.}
\runlabel{Transition paths from the 2023 state: the two Policy Results slides.}
\end{frame}

\begin{frame}{Conclusion}
\begin{itemize}\setlength\itemsep{1.4em}
  \item A lifecycle model of fertility with realistic housing: tenure, down payments, space
  \item Children need space; large homes must be bought; the old hold most of them
  \item Today is a point on a transition to a smaller population
  \item A rebated property tax can move space toward young families, generating a fertility response
\end{itemize}
\end{frame}

\begin{frame}[plain,c,noframenumbering]{}
\centering{\Huge Thank You!}
\end{frame}

\appendix

\begin{frame}[plain,c,noframenumbering]{}
\centering{\Large\textbf{Appendix}}
\end{frame}



\begin{frame}{The 2023 Equilibrium}
\centering
% Display copy of the price and children panels of patch_readout_fit/figures/equilibrium_2023.pdf, drawn from figure_verification.json.
\includegraphics[width=0.9\textwidth,height=0.8\textheight,keepaspectratio]{figures/equilibrium_2023_prices_children.pdf}
\vspace{-0.6em}\hole{Sept 14 point; to be redrawn from the 100-period 2023 state.}
\end{frame}


\begin{frame}{Calibration: Full Fit with Weights}
\scriptsize
\renewcommand{\arraystretch}{0.88}
\input{\tabdirv/fit_full.tex}\par
\smallskip
\scriptsize norm.: normalization, matched exactly; val.: validation, weight zero. Loss contributions sum to the weighted loss \baseloss.
Current base: Mac round 3, chain \basechain; house price \baseprice, $H_0=\baseHzero$.
\end{frame}

\begin{frame}[label=large-homes-data]{Who Owns the Large Homes}
\titlelink{intro-channels}{Back}
\centering
\includegraphics[width=0.9\textwidth,height=0.78\textheight,keepaspectratio]{figures/large_homes_by_age_children_acs2023.pdf}
\par\smallskip
{\scriptsize Owner-occupied homes with $6+$ rooms; ACS 2023.}
\end{frame}

\begin{frame}[label=big-units-owned]{Big Units Are Owned}
\titlelink{housing-constraints}{Back}
\centering
\includegraphics[width=0.77\textwidth,height=0.77\textheight,keepaspectratio]{../code/data/ahs_supply_snapshot/output_ahs_family_unit_menu_national/ahs_tenure_share_within_size.png}
\source{American Housing Survey, 2023.}
\end{frame}

\begin{frame}[label=empirics]{Data}
\titlelink{calibration-overview}{Back}
\begin{itemize}\setlength\itemsep{0.75em}
 \item \textbf{AHS and ACS:} housing sizes, tenure, household composition, and age profiles.
 \item \textbf{PSID:} housing around childbirth, earnings, and wealth by age.
 \item \textbf{CPS and natality records:} completed fertility, childlessness, and birth timing.
 \item \textbf{Census population estimates:} age--sex population and demographic flows.
\end{itemize}

\vfill
Cross-sectional housing patterns discipline access to space; longitudinal data describe housing adjustment around births.
\end{frame}

\begin{frame}{Ownership and Housing Size Around First Birth}
\begin{columns}[T]
\begin{column}{0.48\textwidth}
\centering
{\small Ownership probability}

\smallskip
\includegraphics[width=\linewidth]{JMP_slides/assets/own_f_c_y_all.png}
\end{column}
\begin{column}{0.48\textwidth}
\centering
{\small Number of rooms}

\smallskip
\includegraphics[width=\linewidth]{JMP_slides/assets/rooms_f_c_y_all.png}
\end{column}
\end{columns}

\medskip
\small Original May event-study coefficients; measurement revisions under review.
\source{PSID; original May event-study estimates.}
\end{frame}

\begin{frame}{Moving Around First Birth}
\begin{columns}[T]
\begin{column}{0.48\textwidth}
\centering
{\small Expansion or better housing}

\smallskip
\includegraphics[width=\linewidth]{JMP_slides/assets/mv_s_f_c_y_all.png}
\end{column}
\begin{column}{0.48\textwidth}
\centering
{\small Neighborhood, schools, or proximity}

\smallskip
\includegraphics[width=\linewidth]{JMP_slides/assets/mv_n_f_c_y_all.png}
\end{column}
\end{columns}

\medskip
\small Original May event-study coefficients; measurement revisions under review.
\source{PSID; original May event-study estimates.}
\end{frame}


% CES variant appendix (added Oct 8 2026, review build). Sources:
%   specification and contract: output/model/ces_calibration_20261008/PLAN.md (sec. 1-2), contract_v1.json;
%   matrix: output/model/ces_composite_20261007/tables_matrix.md (cap-6 cells; also cap55_replication_20261008/ces_matrix_cap6_vs_cap55.md);
%   start-2 experiments: docs/model/two_day_synthesis_20261008/sources.md (section 6).
\begin{frame}[label=ces-appendix]{CES Housing Composite}
\titlelink{loan-not-insurance}{Back}
\footnotesize
\setlength{\abovedisplayskip}{3pt}\setlength{\belowdisplayskip}{3pt}
Replace the Cobb--Douglas composite and the parents' space floor $h_P$ with a CES composite in which children raise the rooms needed for given housing services:
\[
u_t=\frac{1}{1-\sigma}\left(\frac{Q_t}{e(m_t)}\right)^{1-\sigma}+\psi_t m_t^{1-\gamma},\qquad
Q_t=\Bigl[\alpha_0\,c_t^{\rho}+(1-\alpha_0)\Bigl(\frac{\chi^{\1\{\mathrm{own}\}}\,h_t}{k(m_t)}\Bigr)^{\rho}\Bigr]^{1/\rho},
\]
\[
\rho=\frac{\varepsilon-1}{\varepsilon},\qquad \varepsilon=0.67,\qquad
k(m_t)=\exp\bigl[\eta_1\1\{m_t\ge1\}+\eta_2\max(m_t-1,0)\bigr].
\]
\begin{itemize}\setlength\itemsep{2pt}
  \item Housing and consumption are complements ($\varepsilon<1$); with no floor, space is not a fixed entry cost of a child.
  \item Estimated: $\eta_1$, $\eta_2$, rooms needed by the first and each later child at home; $\alpha_0$, weight on consumption.
  \item Twelve parameters, thirteen moments. New: rent share of renters 30--55 with no child at home (0.360, CE), renter-to-owner in four years, ages 22--41 (0.281, PSID), mothers with 3+ children (0.357, CPS). Re-measured: rooms added at first birth (1.025, PSID). Rooms with 3+ versus 1--2 children becomes a target; the recent-parent ownership gap becomes validation.
  \end{itemize}
\end{frame}

\begin{frame}[t]{CES Variant: Credit Across Specifications}
\titlelink{loan-not-insurance}{Back}
\footnotesize
Before recalibration: each row re-fits only $\psi$ to childlessness (19\%); house prices fixed, rental cap 6.

\scriptsize\setlength{\tabcolsep}{4pt}\renewcommand{\arraystretch}{0.95}
\begin{tabularx}{\textwidth}{@{}Xccc rr rr@{}}
\toprule
 & & & & \multicolumn{2}{c}{Financed share $80\%\to95\%$} & \multicolumn{2}{c}{Renter credit, 2 years} \\
\cmidrule(lr){5-6}\cmidrule(l){7-8}
 & $\varepsilon$ & Floor & Need $k$ & Impact & Completed & Impact & Completed \\
\midrule
Cobb--Douglas (benchmark) & 1 & yes & 1 & $+0.42\%$ & $-1.35\%$ & $+8.5\%$ & $+0.62\%$ \\
CES & 0.5 & yes & 1 & $-0.16\%$ & $-1.40\%$ & $+8.0\%$ & $+0.22\%$ \\
CES & 0.25 & yes & 1 & $-0.06\%$ & $-0.87\%$ & $+7.3\%$ & $-0.25\%$ \\
CES & 0.5 & yes & 1.5 & $+0.05\%$ & $-1.10\%$ & $+7.0\%$ & $-0.63\%$ \\
CES, no floor & 0.5 & no & 1 & $+0.44\%$ & $-0.03\%$ & $+6.9\%$ & $+2.25\%$ \\
CES, no floor & 0.5 & no & 1.5 & $+0.51\%$ & $-0.18\%$ & $+8.3\%$ & $+2.39\%$ \\
Cobb--Douglas, no floor & 1 & no & 1.5 & $+1.03\%$ & $+0.23\%$ & $+8.3\%$ & $+2.42\%$ \\
\bottomrule
\end{tabularx}
\smallskip

\footnotesize
\begin{itemize}\setlength\itemsep{1pt}\setlength\topsep{2pt}
\item Mortgage credit moves completed fertility by $-2\%$ to $+0.3\%$ in all thirteen specifications that reach 19\% childlessness; the negative sign comes mostly from the floor.
\item Renter credit raises births on impact by 6--9\%; completed fertility, $-0.6\%$ to $+2.4\%$ in the cases shown.
\end{itemize}
\runlabel{Fixed price, rebate held. Impact: births on the baseline distribution of states. Completed: children ever born at 46--49. $k$: rooms-needed factor with a child at home.}
\end{frame}

\begin{frame}[t]{CES Variant: Calibrated Fit}
\titlelink{loan-not-insurance}{Back}
\footnotesize
\begin{center}
\placeholder{[Placeholder: target-fit table (13 moments: target, model, weight) and the 12 parameter estimates,\\
from the overnight multistart calibration, to be filled once it is selected.]}
\end{center}
\vfill
At a starting point of the search (not a fit), completed fertility responds:
\smallskip

\begin{center}\scriptsize
\begin{tabular}{@{}lr@{}}
\toprule
Experiment & Completed fertility \\
\midrule
Financed share $80\%\to100\%$ & $-0.4\%$ \\
Lower earnings risk, mean-preserving & $+13.1\%$ \\
Same, with 100\% financing & $+13.8\%$ \\
\bottomrule
\end{tabular}
\end{center}
\vfill
\hole{Calibration running overnight (12 parameters, 13 moments, $\varepsilon=0.67$); replace the placeholder and the starting-point table with the selected fit and its experiments.}
\end{frame}


\input{appendix_solution_algorithms.tex}

\end{document}
